% =============================================================================
% Discussion
% =============================================================================

\section{Discussion}
\label{sec:discussion}

\subsection{What this means for deployment}

The headline takeaway from \S\ref{sec:evaluation} is that cacheblend's
rescue is real but not free. On Llama-3.3-70B-Instruct fp8 with our
patches, the rescue is bimodal (median 76.5\% across the main
$n{=}99$ corpus, with 6\% of turns reaching the published 95--99\%
peak; Table~\ref{tab:hit-rate}, Table~\ref{tab:cross-corpus}) ---
and at $T{=}0$, severe ($<$0.3 token-identity to no-cache) divergence
appears on 30\% of substantive turns in the deep-evaluation subset
($n_{\text{sub}}{=}47$) and on 26\% in the combined main corpus
($n{=}99$). On the divergent slice of the deep subset
($n_{\text{div}}{=}19$), an LLM judge prefers cacheblend's output
on only 31.6\% of pairs (Table~\ref{tab:judge}). The asymmetry has
three deployment-relevant implications.

\paragraph{For hit-rate-only workloads, the quality cost is irrelevant
--- but the latency cost still applies.} Any deployment that uses
prefix caching for throughput accounting without surfacing model
outputs directly to end users --- bulk embedding pipelines,
throughput-benchmarking harnesses, KV-cache research itself --- need
not worry about the 30\% divergence: those outputs aren't read. But
the TTFT result of \S\ref{sec:eval-latency} (cacheblend 55\% slower at
median, 7$\times$ slower at peak rescue) still applies: cache hits
are not free latency on a 70B fp8 + H100 backend because LOAD
bandwidth approximately cancels prefill savings. Whether cacheblend
is a net win for these workloads depends on the deployment's
LOAD-to-prefill cost ratio (smaller models, faster interconnect,
slower compute, or throughput-vs-latency optimization all shift the
answer).

\paragraph{For interactive coding agents like Claude Code, the
divergence is load-bearing.} The 37\% no-cache-preferred rate is not a
rounding error: an agent that picks a measurably worse next step on
roughly one in three post-cache turns will surface to end users as
``the cached path is dumber.'' Worse, the failure mode is sometimes a
\emph{decode-mode flip} (tool call$\to$prose, 11\% of substantive
turns); a CC harness that expects a tool-call JSON and receives prose
must either retry, surface a degraded user experience, or silently
fall through. None of these are good outcomes.

\paragraph{For batch RAG-style workloads, the published cacheblend
evaluation remains valid.} Our finding is specific to agent traffic
with multi-turn cache state accumulation; RAG-style workloads where
cache LOOKUP fires once against a single retrieved document set do not
exhibit the same dual-mode divergence pattern. The cacheblend paper's
RAG measurements~\cite{lmcache-cacheblend} are not contradicted by ours.

\subsection{Why agent traffic is harder than RAG for cacheblend}
\label{sec:disc-why-agents}

The cacheblend paper~\cite{lmcache-cacheblend} reports near-100\%
output-quality preservation on permuted-RAG benchmarks at the same
0.15 recompute ratio we use. Our 60\% bit-identity + 30\%
severe-divergence + 8.5\% asymmetric mode-flip rate is a much worse
quality picture on the same mechanism. The 0.15 ratio did not
change; the workload did. We hypothesize three structural reasons,
which we are not yet able to fully ablate but state as testable
hypotheses for follow-up work:

\paragraph{(i) Decode mode is binary in agent traffic, continuous in
RAG.} A RAG response is a prose answer where a small logit
perturbation at position $k$ produces a slightly different word at
position~$k$; the rest of the answer remains semantically close. An
agent response in CC's protocol must commit, within the first $\sim$5
output tokens, to one of two modes: emit \texttt{\{"type":~"function"}
(tool-call) or emit natural-language prose (which the harness cannot
act on). Cacheblend's selective recompute introduces a bounded
logit perturbation, but at position~0 of the output that perturbation
can flip the argmax from \texttt{\{} to \texttt{T}/\texttt{B}/\texttt{L}
(common prose openers). Once a mode is chosen, attention to the
already-emitted prefix reinforces it. RAG outputs do not have this
binary regime sensitivity; agent outputs do, and
\S\ref{sec:eval-mode}'s 8.5\% asymmetric stall rate is the direct
consequence.

\paragraph{(ii) Multi-turn cache accumulation amplifies drift on shared
chunks.} RAG is typically single-turn: retrieve, blend, answer. Agent
traffic is multi-turn, and shared chunks between turns are common ---
the system prompt, repeated tool definitions, and prior conversation
history all appear in successive prompts. \emph{In our replay setup,
prompts are fixed captured request bodies; the model's generated
outputs do not feed back into subsequent prompts.} The drift
accumulation is purely \emph{cache-side}: turn $N$'s STORE writes KV
for the prompt chunks it just processed --- including any chunks
whose KV was itself produced by cacheblend's selective recompute on
turn~$N{-}1$. When turn~$N{+}1$ LOOKUPs a chunk that overlaps with
turn~$N$'s prompt, it retrieves whatever KV turn~$N$ wrote, which
may already be one blending step away from fresh compute. The
cacheblend paper's ``small but bounded'' drift bound applies to each
individual blend operation; we hypothesize that on chunks shared
across multiple consecutive turns the cumulative drift compounds.
This is a hypothesis we have not directly isolated --- a clean test
would require comparing single-turn vs multi-turn replays at fixed
$N$th-turn cache states, which we do not run here.

\paragraph{(iii) Templated prompts have positional dependencies that
permuted-RAG does not.} The cacheblend mechanism is most accurate when
the model's attention to a chunk is approximately invariant to that
chunk's exact position --- permuted-RAG, by construction, trains on
this kind of invariance. CC's prompts are highly templated (system
prompt with tool list at fixed offsets, conversation history in
chronological order, current user message at a known position) and the
model has learned strong position-conditioned patterns. When
cacheblend re-uses a cached chunk at a new position, the chunk's KV
reflects attention computed in the \emph{original} positional context,
and the selective-recompute correction sees only 15\% of layers
through the new positional context. For position-sensitive templates
the residual drift on the un-recomputed 85\% of layers compounds; for
permuted-RAG it largely cancels.

Hypothesis~(iii) is directly testable: at
\texttt{LMCACHE\_BLEND\_\allowbreak{}RECOMPUTE\_RATIOS}=1.0, every layer
is recomputed in the new positional context, so the un-recomputed
residual vanishes and identity to no\_cache should recover. We test
this directly in \S\ref{sec:disc-knobs} and find it does. Hypotheses
(i)~and (ii) remain open empirical questions for future work.

\subsection{Recompute-ratio sweep tests hypothesis (iii)}
\label{sec:disc-knobs}

To probe hypothesis~(iii) of \S\ref{sec:disc-why-agents} ---
residual drift on the un-recomputed layers --- we sweep
\texttt{LMCACHE\_BLEND\_\allowbreak{}RECOMPUTE\_RATIOS} over
$\{0.15, 0.5, 1.0\}$ on the
\texttt{swebench\_verified/\allowbreak{}pytest-7490} fixture (4
turns, 3 substantive at $\geq$256 prompt tokens, samples=1, otherwise
identical to the SWE-V+\texttt{post\_compact} subset methodology):

\begin{table}[h]
\small
\centering
\begin{tabular}{lrrr}
\toprule
Ratio & Mean token-identity to NC & Aggregate rescue \% \\
\midrule
\textbf{0.15 (default)} & 0.847 & 27.4\% \\
0.50                    & 0.831 & 27.4\% \\
\textbf{1.00}           & \textbf{1.000} & 27.4\% \\
\bottomrule
\end{tabular}
\caption{Recompute-ratio sweep on \texttt{pytest-7490} ($n{=}3$
substantive turns per ratio). Identity to the no\_cache baseline
recovers to 1.000 at ratio=1.0 (every layer recomputed in the new
positional context) but stays around 0.83 at the intermediate ratio,
not 0.92 as one would extrapolate. Rescue percent is unchanged because
\texttt{RECOMPUTE\_RATIOS} controls layer recompute, not which KV
tokens are loaded from cache.}
\label{tab:recompute-sweep}
\end{table}

Two findings:

\paragraph{Hypothesis (iii) is supported.} At ratio=1.0, identity to
\texttt{no\_cache} recovers bit-perfectly. Recomputing every layer in
the new positional context eliminates the residual drift that the
default 0.15 ratio carries. The mechanism we hypothesized in
\S\ref{sec:disc-why-agents}(iii) is the load-bearing one for the
identity gap.

\paragraph{Partial recompute did not partially recover on this
probe.} Going from $0.15 \to 0.5$ (3.3$\times$ more layers recomputed)
barely moves identity (0.847~$\to$~0.831), within noise on $n{=}3$.
Recovery is concentrated at the discrete $\to{}1.0$ jump on the data
points we have. \emph{We do not claim the intermediate-ratio tradeoff
curve is established from $n{=}3$:} a larger sweep (denser ratio grid
$\times$ the full deep-evaluation subset) is needed before
concluding that no monotonic recovery exists between 0.15 and 1.0.
What is robust from this probe is the discrete endpoint contrast:
ratio~1.0 produces identity~$=$~1.000 by construction (every layer
recomputed in the new positional context), confirming the
positional-context hypothesis. Whether 0.5 or 0.8 would offer a
useful Pareto-improvement remains an open question.

The combination of \S\ref{sec:eval-latency}'s TTFT measurement
(cacheblend already slower than \texttt{no\_cache} at the default
0.15 ratio on our backend) and this sweep's discrete-jump pattern
suggests cacheblend's default ratio is not clearly a good operating
point in our measurements, and that ratio=1.0 may not be the
practical fix (it recomputes every layer, eliminating the rescue
benefit). A genuine graded tradeoff curve --- one that gives
proportional identity recovery against linear cost --- may require
mechanism changes beyond the existing recompute-ratio knob
(smaller cacheable blocks, layer-importance-aware recompute
selection, or model-architecture co-design are candidates worth
exploring).

\paragraph{Sample size caveat.} Three substantive turns per ratio is
a small sample; the 0.847~$\to$~0.831 dip at 0.5 is plausibly noise.
What is robust is the discrete jump at 1.0 from $\sim$0.84 to exactly
1.000 --- that comes from the math, not the sample. A larger probe
across the full SWE-V+\texttt{post\_compact} subset would tighten the intermediate-ratio
identity estimates but is unlikely to change the discrete
0.15$\leftrightarrow$1.0 endpoint contrast.

\subsection{Resource-constrained methodology choices}
\label{sec:disc-resource}

This paper was produced on a fixed compute budget (\$25 envelope for
the headline experiments, against a 2$\times$ H100 80GB SXM pool with
real per-pod cost). Several methodological tightenings a reviewer
would reasonably ask for --- a raw-no-cache baseline that bypasses our
parser to validate separator injection is content-neutral, an
intra-condition repeatability study to bound run-to-run
nondeterminism, instrumented per-phase TTFT breakdowns (LMCache
LOAD vs prefill vs blend vs scheduler vs decode), a denser
recompute-ratio grid, a quantitative patch ablation that runs each
intermediate configuration as its own pod replay, and a bf16-vs-fp8
precision comparison --- are individually small but cumulatively a substantial increase in
pod-hours, well beyond this project's compute budget. We list them
explicitly as open empirical questions rather than implicit ones, so
that a reader who wants to commission this work knows where it sits
relative to our existing claims. The headline findings (rescue is
bimodal, identity gap is real, latency cost is real, mode-flip
asymmetry is real) are robust to the methodology gaps we name here ---
they are direct read-offs from the measurements we did run, not
inferences that depend on the missing instrumentation. The
\emph{causal} explanations for those findings are where the missing
instrumentation would matter most; we have flagged each one as a
hypothesis rather than a conclusion.

\subsection{Limitations}
\label{sec:disc-limits}

\textbf{Single model.} All measurements are on Llama-3.3-70B-Instruct
fp8. The cacheblend mechanism is model-agnostic in design, but the
empirical quality cost is not: the fp8 quantization noise floor, the
70B parameter count, and Llama's attention architecture all interact
with cacheblend's selective recompute. The same patches and methodology
on Sonnet, Haiku, Llama-8B, or Qwen could produce a different
divergence distribution. \S\ref{sec:disc-knobs}'s recompute-ratio probe
is the higher-priority follow-up; multi-backend characterization
is the longer-horizon next step.

\textbf{Judge calibration.} We replicate the 19-pair A/B against
Opus-4.7 and Haiku-4.5 in \S\ref{sec:eval-judge}, and the
three-judge majority vote reproduces Sonnet-4.6's headline numbers
(31.6\% preference, 63.2\% non-inferior). The single-judge risk is
reduced but not eliminated --- all three judges share Anthropic's
RLHF training, so a coordinated calibration bias could still apply.
A human spot-check is a sensible next step we do not run here.

\textbf{Single decoding temperature.} $T=0$ greedy decoding makes
single-token differences load-bearing: a divergence at position $k$
permanently forks the trajectory. Sampling at $T=0.7$ would average
over multiple plausible continuations and might mask or amplify the
cross-condition divergence in the marginal distribution. We defer
the $T>0$ distributional measurement to follow-up work.

\textbf{Sample size.} 19 divergent pairs is small. The 95\% Wilson CI
on the cacheblend preference rate spans [15.4\%, 54.0\%] --- we cannot
reject ``cacheblend is preferred half the time'' from the data alone.
The point estimate's location well below 50\% is what supports the
non-inferiority gate failing; widening the corpus would tighten the
interval but is unlikely to change the qualitative direction.

\textbf{Workload diversity within \dataset{}, but a single corpus
source.} The combined $n{=}99$ corpus mixes SWE-bench-Verified
sessions, \texttt{post\_compact/} canonical captures, and
\texttt{skill\_invocation/} (\S\ref{sec:eval-cross-corpus}). All three
subsets come from one researcher's laptop running one client family
(CC SDK), so the captured prompt distribution is not yet broad across
operators, languages, or task domains. The cross-corpus replication
of the bimodal-rescue and severe-divergence patterns supports
generality \emph{within Claude-Code-shaped traffic}; we make no claim
about other agent harnesses (Cursor, Aider, OpenAI Codex, etc.) on
this evidence alone.

\textbf{Backend specificity.} All measurements use a
Llama-3.3-70B-Instruct fp8 backend on 2$\times$ H100 80GB SXM, vLLM
v0.19.0, and lmcache v0.4.2 with our patches. We do not claim
generality across smaller models, different precisions (bf16),
faster interconnects, or non-vLLM serving stacks. The bf16 vs fp8
comparison is the highest-priority open empirical question;
\S\ref{sec:disc-knobs}'s recompute sweep already isolates one
hypothesis on the present backend.

\subsection{Recommendation}
\label{sec:disc-recommend}

We make one concrete methodological recommendation:
\emph{KV-cache reuse research should report output-quality measurements
alongside hit-rate measurements, on at least one agentic workload, with
a structured methodology that includes $T=0$ token-identity and a
strong-judge blind A/B preference test.} The hit-rate-only framing
that the field has converged on is sufficient for the RAG-style
workloads the early prefix-caching literature targeted, but it
silently elides exactly the failure mode we surface here: a mechanism
can hit at the published rate and still produce systematically worse
output on the same input. We provide \dataset{} and the \sys{} harness
as the reproducible substrate for this kind of evaluation.
